\honest{Searching for Bayes-Structure}{Eliezer Yudkowsky}{2008}

I open with a quotation about gnomish helms, which should not work and are all bells and whistles. Those that do work have a small working helm hidden inside, disguised as inessential.

We have seen that knowing about something means sharing mutual information with it, and that such information is physical. Forming true beliefs without evidence is ``the same sort of improbability'' as warm water turning into ice and electricity. Rationality takes thermodynamic work, imperfect minds give off heat, and that work is governed by probability theory, of which thermodynamics is a special case. ``Statistical mechanics is a special case of statistics.'' \nb{That is E.~T. Jaynes's view, and it is contested; the note on ``The Second Law of Thermodynamics, and Engines of Cognition'' gives the details. Jaynes is not named here.}

A wheel that turns with no visible power source must have a hidden battery. Likewise a mind that arrives at true beliefs must be doing something ``at least vaguely Bayesian'' somewhere. If at first it shares no information with some part of the world and later shares 10 bits, it must in between have met evidence, which is the same thing as mutual information, and processed some of it in the right direction. The alternative is to have created information from nothing. So any part of a mind that helps find truth ``must have at least a little Bayesian structure.'' \nb{The post does not define ``Bayesian structure.'' So qualified, the claim says little more than that a useful part helps the mind's states carry information about the world, which is the premise.}

Philosophers spent centuries on the meaning of words, and ``all along, it was a disguised form of Bayesian inference!'' I name no philosopher. I was a little disappointed that no one jumped up and cried that it was Bayes all along. Perhaps it is less exciting when someone else unravels the mystery (Newton had more fun than calculus students), or when you do not see how great the quest is.

Searching for hidden Bayesian structure is a quest like the search for the Holy Grail. It is a different quest for each part of the mind, ``but the Grail always turns out to be the same,'' the ``entire'' Grail, found by seeking a full answer in whatever form, not by hand-waving Grailish arguments. Readers told me my long essays did not make clear where I was going, but it is hard to announce a destination like that. My evidence is my own experience: ``Once this happens to you a few times, you kinda pick up the rhythm.'' \nb{The argument about the second law supports only ``at least a little'' Bayesian structure. The claim that the full structure is always found rests on this reported experience.}

I describe the experience. You learn about some useful mental process that philosophers have argued over for centuries and AI researchers cannot agree how to build. It does not look Bayesian, but underneath it is. Then a completely different process turns out Bayesian too. Then I ask whether the parts that still look non-Bayesian do anything. If yes, ``those are Bayesian too!'' If no, ``what a stupid design.'' I could eat a bucket of amino acids and vomit a better brain. Talking about this rhythm is like dancing about architecture. \nb{Both answers confirm the thesis. A useful part without Bayesian structure, the one finding that would count against it, is not among the outcomes.}

I grant that knowing that a process is Bayesian, without knowing how, is ``a hint at the form an answer would take; but certainly not an answer.'' \nb{By that standard the thermodynamic argument gives a hint. It shows that some Bayesian structure exists, not what it is.} If I say ``Bayes is the secret of the universe,'' some will cheer and others will snort at my narrowness and point to useful ad hoc methods such as regularized linear regression. My link shows its Bayesian reading. I hoped that one worked example, plus the distinctions between passwords and knowledge and between tools and laws, would convey some of the rhythm. The full secret is known only to the Bayes Council, ``and if I told you, I'd have to hire you.'' \nb{The link is apt: ridge regression gives the most probable solution under a normal prior.}

The aim is ``Bayes-Sight'': to see through the surface to the probability flows, to know how each process is Bayesian, ``as it always is - as it always must be.'' I close with Zelazny: a queen who might see ``The clear, cold lines of eternity, I daresay. Beneath all Shadow.''
